% !PRINTED none
\weberpage{none}
\begin{center}
\gesperrt{Fünfter Abschnitt.}

\gesperrt{Constitution der allgemeinen Gruppen.}
\end{center}

\rule{4cm}{0.4pt}

\begin{center}
\S.\,27.

\gesperrt{Bildung von Gruppen nach Cayley.}
\end{center}

Die allgemeine Definition der Gruppe, die im \S.\,1 gegeben ist, lässt über die Natur dieses Begriffes noch manches im Dunkel, und auch die verschiedenen einzelnen Gruppen, die wir im Verlauf unserer Betrachtungen kennen gelernt haben, geben nur Hinweise auf allgemeine Gesetze, und zeigen, dass der Gruppenbegriff an sich nichts Widersprechendes hat. In der Definition der Gruppe ist mehr enthalten, als es auf den ersten Blick den Anschein hat, und die Zahl der möglichen Gruppen, die aus einer gegebenen Anzahl von Elementen zusammengesetzt werden können, ist eine sehr beschränkte. Die allgemeinen Gesetze, die hier herrschen, sind erst zum kleinsten Theile erkannt, so dass jede neue specielle Gruppe, namentlich bei kleinerer Gliederzahl, ein neues Interesse bietet und zu eingehendem Studium auffordert.

\gesperrt{Welche Gruppen sind zwischen einer gegebenen Zahl von Elementen, d.\,h. bei gegebenem Grade überhaupt möglich?} Das ist die allgemeine Frage, um die es sich handelt, von deren vollständiger Lösung wir aber noch weit entfernt sind. \gesperrt{Cayley} hat diese Aufgabe zuerst für die niedrigsten Gradzahlen in Angriff genommen\footnote{On the theory of groups, as depending on the symbolic equation $\Theta^n = 1$. Philosophical Magazine, Vol.\ II, 1854. (\gesperrt{Cayley's} mathematical papers, Vol.\ II, 125.) American Journal of mathematics, Vol.\ I.}.
